\subsection{Sobolev Spaces}

Let U be an open subset of \(\mathbf{R}^{n} \). Let p be a real number \(\geqslant 1\) and let k be a natural number. We denote by \(W^{k,p}(U)\) the space of distributions \(f \in \mathcal{D}'(U)\) such that, for every \(\alpha \in N^{n}\) with \(|\alpha| \leqslant k\) , the distribution \(\partial^{\alpha}f\) is associated with an element of \(\mathrm{L}^{p}(U) \).

In particular, when \(U = \mathbf{R}^{n}\) , the elements of \(W^{k,p}(U)\) are tempered distributions.

The map from \(\mathrm{W}^{k,p}(\mathrm{U})\) into \(\mathbf{R}_{+}\) that assigns to \(f\in \mathrm{W}^{k,p}(\mathrm{U})\)

\[
\| f \| _ {k, p} = \Bigl (\sum_ {| \alpha | \leqslant k} \| \partial^ {\alpha} f \| _ {p} ^ {p} \Bigr) ^ {1 / p}
\]

is a norm on \(W^{k,p}(U)\) . The space \(W^{k,p}(U)\) will always be equipped with this norm; this normed space is called the Sobolev space of index k and exponent p.

The space \(\mathcal{D}(\mathrm{U})\) is contained in \(\mathrm{W}^{k,p}(\mathrm{U})\) . We denote \(\mathrm{W}_0^{k,p}(\mathrm{U})\) the closure of \(\mathcal{D}(\mathrm{U})\) in \(\mathrm{W}^{k,p}(\mathrm{U})\) . This is a closed subspace of \(\mathrm{W}^{k,p}(\mathrm{U})\) .

One has \(\mathrm{W}_{0}^{k,p}(\mathbf{R}^{n})=\mathrm{W}^{k,p}(\mathbf{R}^{n})\) , but the spaces \(\mathrm{W}^{k,p}(\mathrm{U})\) and \(\mathrm{W}_{0}^{k,p}(\mathrm{U})\) are distinct in general (cf. exercises 12 of IV, p. 334 and 14 of IV, p. 334).

We also denote \(\mathrm{H}^k (\mathrm{U}) = \mathrm{W}^{k,2}(\mathrm{U})\) and \(\mathrm{H}_0^k (\mathrm{U}) = \mathrm{W}_0^{k,2}(\mathrm{U})\) .

The norm of \(\mathrm{H}^{k}(\mathrm{U})\) is a pre-Hilbertian norm, associated with the positive Hermitian form on \(\mathrm{H}^{k}(\mathrm{U})\) defined by

\[
\left(f _ {1}, f _ {2}\right) \mapsto \sum_ {| \alpha | \leqslant k} \int_ {\mathrm{U}} \overline {{\partial^ {\alpha} f _ {1}}} \partial^ {\alpha} f _ {2} d \mu .
\]

The Hilbert space \(\mathrm{H}^k (\mathrm{U})\) coincides with the space denoted \(\mathcal{H}^k\) in EVT, V, p. 6, example (3).

One has \(\mathrm{W}^{0,p}(\mathrm{U})=\mathrm{L}^{p}(\mathrm{U})\) and \(\mathrm{H}^{0}(\mathrm{U})=\mathrm{L}^{2}(\mathrm{U})\) by definition; moreover \(\mathrm{W}_{0}^{0,p}(\mathrm{U})=\mathrm{L}^{p}(\mathrm{U})\) by prop. 4, b) of IV, p. 202.

PROPOSITION 20. — The Sobolev spaces \(W^{k,p}(U)\) and \(W_{0}^{k,p}(U)\) are Banach spaces of countable type. In particular, the spaces \(H^{k}(U)\) and \(H_{0}^{k}(U)\) are Hilbert spaces of countable type.

It suffices to prove the assertions concerning \(\mathrm{W}^{k,p}(\mathrm{U})\) .

Let I be the set of \(\alpha \in \mathbf{N}^n\) such that \(|\alpha| \leqslant k\) . The linear map \(u\) from \(\mathrm{W}^{k,p}(\mathrm{U})\) into \(\mathrm{L}^p (\mathrm{U})^{\mathrm{I}}\) that associates with \(f\) the family \((\partial^{\alpha}f)_{\alpha \in \mathrm{I}}\) is injective; it is continuous and strict by definition of the norm on \(\mathrm{W}^{k,p}(\mathrm{U})\) . To prove that \(\mathrm{W}^{k,p}(\mathrm{U})\) is complete, it suffices to prove that its image under \(u\) is closed. Now, let \((f_n)_{n \in \mathbf{N}}\) be a sequence in \(\mathrm{W}^{k,p}(\mathrm{U})\) such that \((u(f_n))_{n \in \mathbf{N}}\) converges. Let \((g_{\alpha})_{\alpha \in \mathrm{I}} \in \mathrm{L}^p (\mathrm{U})^{\mathrm{I}}\) be its limit. For \(\alpha \in \mathrm{I}\) , the sequence \((\partial^{\alpha}f_n)_{n \in \mathbf{N}}\) converges in \(\mathrm{L}^p (\mathrm{U})\) to \(g_{\alpha}\) . A fortiori, the convergence takes place in \(\mathcal{D}'(\mathrm{U})\) . Set \(f = g_0\) . For every \(\varphi \in \mathcal{D}(\mathrm{U})\) , it follows that

\[
\begin{array}{r l} & {\langle \partial^ {\alpha} f, \varphi \rangle = (- 1) ^ {| \alpha |} \langle f, \partial^ {\alpha} \varphi \rangle} \\ & {\qquad = \lim _ {n \to + \infty} (- 1) ^ {| \alpha |} \langle f _ {n}, \partial^ {\alpha} \varphi \rangle = \lim _ {n \to + \infty} \langle \partial^ {\alpha} f _ {n}, \varphi \rangle = \langle g _ {\alpha}, \varphi \rangle .} \end{array}
\]

This proves that \(g_{\alpha} = \partial^{\alpha} f\) for all \(\alpha \in I\) , hence \((g_{\alpha})_{\alpha \in I} = u(f)\) belongs to the image of u.

The space \(\mathrm{W}^{k,p}(\mathrm{U})\) is identified by u with a subspace of the space \(\mathrm{L}^{p}(\mathrm{U})^{\mathrm{I}}\) ; the latter is of countable type (prop. 2 of IV, p. 180 and TG, IX, p. 19, cor., (ii))), hence so is \(\mathrm{W}^{k,p}(\mathrm{U})\) (loc. cit., (i)).

PROPOSITION 21. --- Let N be the Euclidean norm on \(\mathbf{R}^{n}\) . Let \(k\) be an integer \(\geqslant 0\) . The Sobolev space \(\mathrm{H}^{k}(\mathbf{R}^{n})\) is the space of \(f \in \mathcal{S}'(\mathbf{R}^{n})\) such that \((1 + \mathrm{N}^{k})\mathcal{F}(f)\) belongs to \(\mathrm{L}^{2}(\mathbf{R}^{n})\) .

We proceed by induction on \(k\) . When \(k = 0\) , the result is a consequence of the Plancherel theorem (II, p. 215, th. 1). Suppose that \(k = 1\) . For \(f \in \mathcal{S}'(\mathbf{R}^n)\) , we have \((1 + \mathrm{N})\mathcal{F}f \in \mathrm{L}^2(\mathbf{R}^n)\) if and only if \(f \in \mathrm{L}^2(\mathbf{R}^n)\) and \(\mathrm{N}\mathcal{F}f \in \mathrm{L}^2(\mathbf{R}^n)\) . Moreover, \(\mathrm{N}\mathcal{F} \in \mathrm{L}^2(\mathbf{R}^n)\) if and only if, for every \(\alpha \in \mathbf{N}^n\) such that \(|\alpha| = 1\) , we have \(\mathrm{X}^\alpha\mathcal{F}f \in \mathrm{L}^2(\mathbf{R}^n)\) . Since one has \(\mathcal{F}(\partial^\alpha f) = 2i\pi\mathrm{X}^\alpha\mathcal{F}f\) , this condition means that \(\mathcal{F}(\partial^\alpha f) \in \mathrm{L}^2(\mathbf{R}^n)\) for every \(\alpha\) with \(|\alpha| = 1\) , that is, \(\partial^\alpha f \in \mathrm{L}^2(\mathbf{R}^n)\) for every \(\alpha\) with \(|\alpha| = 1\) . It follows that the assertion is true for \(k = 1\) .

Suppose now that \(k \geqslant 2\) and that the assertion concerning \(\mathrm{H}^{\ell}(\mathbf{R}^{n})\) is valid for every positive integer \(\ell \leqslant k - 1\) . Let \(f \in \mathcal{S}'(\mathbf{R}^{n})\) . By definition, one has \(f \in \mathrm{H}^{k}(\mathbf{R}^{n})\) if and only if \(f \in \mathrm{L}^{2}(\mathbf{R}^{n})\) and, for every \(\beta \in \mathbf{N}^{n}\) such that \(|\beta| \leqslant 1\) , the distribution \(\partial^{\beta} f\) belongs to \(\mathrm{H}^{k-1}(\mathbf{R}^{n})\) . This is equivalent, by the induction hypothesis, to \(f \in \mathrm{L}^{2}(\mathbf{R}^{n})\) and \((1 + \mathrm{N}^{k-1})\mathcal{F}(\partial^{\beta} f) \in \mathrm{L}^{2}(\mathbf{R}^{n})\) for every \(\beta \in \mathbf{N}^{n}\) such that \(|\beta| \leqslant 1\) . Since \(\mathcal{F}(\partial^{\beta} f) = (2i\pi X)^{|\beta|}\mathcal{F}(f)\) , the condition \(f \in \mathrm{H}^{k}(\mathbf{R}^{n})\) is equivalent to saying that \(\mathcal{F}f \in \mathrm{L}^{2}(\mathbf{R}^{n})\) and \((1 + \mathrm{N}^{k-1})\mathrm{X}^{\beta}\mathcal{F}f \in \mathrm{L}^{2}(\mathbf{R}^{n})\) for every \(\beta \in \mathbf{N}^{n}\) such that \(|\beta| \leqslant 1\) .

The inequalities

\[
1 + \mathrm{N}^{k}\leqslant 1 + \mathrm{N}^{k - 1}\sum_{\substack{\beta \in \mathbf{N}^{n}\\ |\beta |\leqslant 1}}|\mathrm{X}^{\beta}|\leqslant 1 + n^{1 / 2}\mathrm{N}^{k}\leqslant (1 + n^{1 / 2})(1 + \mathrm{N}^{k})
\]

then imply that \(f \in \mathrm{H}^k(\mathbf{R}^n)\) if and only if \((1 + \mathrm{N}^k)\mathcal{F} \in \mathrm{L}^2(\mathbf{R}^n)\) .

